% ===========================================================
% Введение ВКР — LaTeX (Overleaf)
% Подключить в основном документе: \input{introduction.tex}
% Убедитесь, что в преамбуле указано:
%   \usepackage[utf8]{inputenc}
%   \usepackage[T2A]{fontenc}
%   \usepackage[russian]{babel}
%   \bibliographystyle{gost-numeric}  % или другой ГОСТ-стиль
%   \bibliography{references}
% ===========================================================

\section*{ВВЕДЕНИЕ}
\addcontentsline{toc}{section}{Введение}

% В условиях глобализации деловых коммуникаций задача точной и безопасной расшифровки многоязычных диалогов приобретает критическое значение для широкого спектра профессиональных сфер. Международные переговоры, судебные заседания, пресс-конференции и рабочие совещания с участием представителей различных языковых групп формируют постоянно растущий объём аудиоданных, требующих структурированной текстовой фиксации с разделением по спикерам и сохранением исходных языков высказываний.

% Существующие технологические решения в области автоматического распознавания речи (Automatic Speech Recognition, ASR) не удовлетворяют в полной мере требованиям пользователей, формируя ряд фундаментальных проблем. Во-первых, наблюдается технологическая фрагментарность: специалисты вынуждены самостоятельно интегрировать разрозненные инструменты для диаризации (разделения речи по спикерам), определения языка (Language Identification, LID) и непосредственно транскрибации. Такой процесс требует глубоких технических компетенций и не гарантирует стабильного результата, особенно при работе с редкими языками или речью с выраженным акцентом.

% Во-вторых, использование облачных сервисов распознавания речи сопряжено с рисками для конфиденциальности данных. Передача аудиозаписей переговоров, содержащих коммерческую, юридическую или журналистскую тайну, на серверы третьих сторон ставит под угрозу безопасность информации и может нарушать требования законодательства о защите персональных данных и внутренних политик организаций.

% В-третьих, существующие решения для локального развертывания либо ограничены предустановленными моделями и не позволяют адаптироваться под специфические задачи пользователя, либо требуют дорогостоящего специализированного оборудования с графическими ускорителями (GPU), что делает их недоступными для большинства потенциальных пользователей. Таким образом, на рынке отсутствует универсальное, конфиденциальное и технологически доступное решение, объединяющее функции диаризации, определения языка и транскрибации в рамках единой расширяемой системы, способной функционировать на стандартном пользовательском оборудовании.

% Указанные обстоятельства обусловливают актуальность разработки модульной программной системы, обеспечивающей полный цикл локальной обработки многоязычных аудиозаписей с использованием pipeline-архитектуры и поддержкой подключаемых моделей машинного обучения.

% Проблематика автоматического распознавания речи активно исследуется в работах отечественных и зарубежных учёных. Значительный вклад в развитие методов ASR внесли исследования, связанные с моделью Whisper~\cite{radford2023whisper}, архитектурой Conformer~\cite{gulati2020conformer} и подходами на основе CTC~\cite{graves2006ctc} и механизма внимания~\cite{chan2016las, vaswani2017attention}. Современные модели, такие как wav2vec~2.0~\cite{baevski2020wav2vec}, продемонстрировали возможности самообучения на неразмеченных данных. Задачи диаризации спикеров рассмотрены в работах H.~Bredin и др.~\cite{bredin2023pyannote, bredin2020pyannote, plaquet2023powerset}, а методы идентификации языка по аудиоданным~--- в исследованиях, посвящённых архитектуре ECAPA-TDNN~\cite{desplanques2020ecapa} и набору данных VoxLingua107~\cite{valk2021voxlingua107}. Вместе с тем вопросы проектирования интегрированных модульных систем, объединяющих указанные компоненты в единый pipeline с поддержкой расширяемости через plugin-архитектуру~\cite{gamma1995patterns, martin2017clean} и ориентированных на локальное развертывание без GPU, остаются недостаточно проработанными.

% \textbf{Цель работы}~--- проектирование и реализация модульной программной системы оффлайн-транскрибации многоязычной речи с использованием pipeline-архитектуры, обеспечивающей полный цикл обработки аудиозаписей встреч (диаризация, определение языка, транскрибация) в локальном режиме на стандартном пользовательском оборудовании.

% Для достижения поставленной цели необходимо решить следующие \textbf{задачи}:
% \begin{enumerate}
%     \item провести анализ существующих методов и моделей автоматического распознавания речи~\cite{radford2023whisper, gulati2020conformer, baevski2020wav2vec}, диаризации спикеров~\cite{bredin2023pyannote} и определения языка~\cite{valk2021voxlingua107, desplanques2020ecapa}, а также выполнить сравнительный обзор существующих программных решений в данной области;
%     \item сформулировать функциональные и нефункциональные требования к разрабатываемой системе;
%     \item спроектировать pipeline-архитектуру обработки аудиоданных~\cite{hohpe2003enterprise, fowler2002patterns}, обеспечивающую последовательное выполнение этапов диаризации, сегментации, определения языка и транскрибации;
%     \item разработать plugin-систему~\cite{gamma1995patterns} для подключения и замены моделей распознавания речи без модификации основного кода системы;
%     \item реализовать механизмы воспроизводимости результатов обработки, включая сохранение конфигурации запусков, промежуточных данных и возможность возобновления прерванного pipeline;
%     \item реализовать графический интерфейс пользователя на базе PyQt6~\cite{pyqt6}, обеспечивающий загрузку аудиофайлов, запуск обработки и отображение результатов;
%     \item провести экспериментальную оценку качества работы системы на многоязычных аудиозаписях с использованием метрик WER~\cite{morris2004wer} и DER~\cite{park2022review_der}.
% \end{enumerate}

% \textbf{Объектом исследования} является процесс автоматической обработки многоязычных аудиозаписей с разделением речи по спикерам и определением языка высказываний.

% \textbf{Предметом исследования} выступает программная архитектура модульной системы оффлайн-транскрибации, построенная на основе pipeline-обработки с поддержкой подключаемых моделей машинного обучения.

% \textbf{Методы исследования.} В работе используются методы системного анализа и сравнительного анализа программных решений, объектно-ориентированного проектирования~\cite{gamma1995patterns}, а также паттерны архитектурного проектирования (Pipeline~\cite{hohpe2003enterprise}, Strategy, Plugin/Registry, Observer). Для оценки качества системы применяются методы экспериментального тестирования с использованием стандартных метрик оценки качества распознавания речи (WER)~\cite{morris2004wer} и диаризации (DER)~\cite{park2022review_der}. Тестирование проводится на открытых бенчмарках~\cite{panayotov2015librispeech, conneau2022fleurs}.

% \textbf{Практическая значимость} работы определяется тем, что разработанная система представляет собой готовое кроссплатформенное программное приложение, пригодное для использования в организациях, предъявляющих повышенные требования к конфиденциальности обрабатываемых данных. Система может применяться для протоколирования международных переговоров и совещаний, фиксации судебных заседаний и допросов, расшифровки интервью и пресс-конференций. Модульная архитектура обеспечивает возможность адаптации системы под специфические задачи пользователей путём подключения специализированных моделей распознавания речи.

% \textbf{Структура работы.} Выпускная квалификационная работа состоит из введения, трёх глав, заключения, списка литературы и приложений.

% В первой главе проводится анализ предметной области: формализуется задача транскрибации многоязычной речи, выполняется обзор существующих моделей ASR~\cite{radford2023whisper, gulati2020conformer}, диаризации~\cite{bredin2023pyannote} и определения языка~\cite{valk2021voxlingua107}, а также сравнительный анализ существующих программных решений. По результатам анализа формулируются требования к разрабатываемой системе.

% Во второй главе описывается проектирование архитектуры системы: общая компонентная архитектура~\cite{martin2017clean}, pipeline-механизм обработки аудиоданных~\cite{hohpe2003enterprise}, модель данных, plugin-система и реестр моделей~\cite{gamma1995patterns}, механизмы воспроизводимости и отказоустойчивости, событийная модель взаимодействия с графическим интерфейсом.

% В третьей главе представлены результаты реализации системы: описаны используемые технологии и инструменты~\cite{ravanelli2021speechbrain, faster_whisper, pytorch, pyqt6}, ключевые аспекты реализации pipeline-ядра, plugin-системы и графического интерфейса. Приводятся результаты экспериментальной оценки качества работы системы на многоязычных аудиозаписях.